\origpage{76}{ROBERTSON AND SEYMOUR}
\noindent (!) If $i \le m$ and $|V(T_i)| = 1$ then $|E(B_i)| = 1$ and $|V(B_i)| \le 2$. We determine the folio of $Z_i$ in $B_i$ and return to (1) for the next iteration.

\hspace*{12pt}(2) If $i \le m$ and $|V(T_i)| > 1$, let $f_j, f_k$ be the two edges of $T_i$ with a common end with $f_i$. Then $j, k < i$, and $(B_j, B_k)$ is a separation of $B_i$, and the $\delta$-folios of $B_j$ relative to $Z_j$ and $B_k$ relative to $Z_k$ have already been determined. Since $V(B_j \cap B_k) \subseteq Z_j \cap Z_k$ we can calculate the $\delta$-folio of $B_j \cup B_k$ relative to $Z_j \cup Z_k$ from this information, by (2.2). Since $Z_i \subseteq Z_j \cup Z_k$ and $B_j \cup B_k = B_i$, the $\delta$-folio of $B_i$ relative to $Z_i$ can be calculated by (2.1). (Note that, since $\delta$-folios have size bounded by a constant, these calculations take constant time.) We return to (1) for the next iteration.

\hspace*{12pt}(3) If $i = m + 1$, the folio of $B_m$ relative to $Z_m$ has been determined. We determine the $\delta$-folio of $A_m$ relative to $(Z \cup V(B_m)) \cap V(A_m)$ (this is easy, since $|V(A_m)| \le 2$), and we use (2.1) and (2.2) to determine the $\delta$-folio of $Z$ in $A_m \cup B_m = G$.

\hspace*{12pt}(4) For each member of the $\delta$-folio of $Z$ in $G$ (up to isomorphism) we backtrack through this computation to compute a corresponding model, in the natural way.

Each iteration takes constant time, and there are $|E(T)| + 1$ iterations. Thus, steps (1)-(3) take time $O(\varepsilon)$. In step (4), the backtracking to determine how each edge of $G$ should be used to produce the desired member of the $\delta$-folio takes time $O(\varepsilon)$ and converting this information into a model again takes time $O(\varepsilon)$. Thus the total running time is $O(\varepsilon)$. $\square$

To apply (4.1), we use the following theorem, a consequence of Theorems (1.5) of [15] and (5.1) of [19]. See also [25] for a numerical improvement.

\stmt
(4.2) \textsc{Theorem}. \textit{For every planar graph $H$ there is a number $\omega$ such that every graph with branch-width $\ge \omega$ has a minor isomorphic to $H$.}

For a first application of (4.1), we have the following, improving a result of [13].

\stmt
(4.3) \textit{Let $\mathcal{F}$ be a minor-closed class of graphs such that some planar graph is not in $\mathcal{F}$. Then there is an algorithm to determine whether $G \in \mathcal{F}$ with running time $O(v^2)$.}

\textit{Proof.} By Theorem (1.4) of [14], there is a finite list of graphs $H_1, \dots, H_k$ such that a graph $G$ is in $\mathcal{F}$ if and only if it has no minor isomorphic to any of $H_1, \dots, H_k$. Since $\mathcal{F}$ does not contain every planar graph, some $H_i$ (say $H_1$) is planar. By (4.2) there is a number $\omega \ge 2$ such that every graph with branch-width $\ge \omega$ has a minor isomorphic to $H_1$. Let $\delta = \max(\varepsilon(H_1), \dots, \varepsilon(H_k))$. We now describe an algorithm to test membership of $\mathcal{F}$. Let $G$ be the input graph. We may assume that $|E(G)| \ge 2$.
